%!TEX root = ../note.tex
\section{Introduction}
Calculus has two branches. \emph{Differential calculus} studies slopes and
tangent lines; \emph{integral calculus} studies areas and accumulation. Both
rest on the notion of a limit, and limits only behave well because of a
property of the real numbers that the rationals lack. This course therefore
starts by pinning down exactly what the real numbers are.

We first recall the number systems $\N\subseteq\Z\subseteq\Q$ and the language
of statements and quantifiers in which everything else is written, together
with induction and the well-ordering principle, the two basic tools for
proving statements about all natural numbers.

We then describe $\R$ as a \emph{complete ordered field}. The field axioms
let us compute, the order axioms let us compare, and the completeness axiom
--- every nonempty set that is bounded above has a least upper bound --- is
what separates $\R$ from $\Q$. Most of the first chapter is about what
completeness buys us: the Archimedean property, the existence of $\sqrt2$,
and the density of $\Q$ and of the irrationals in $\R$.

Finally, we define what it means for a sequence of real numbers to converge.
This is the first genuine limit in the course, and the $\epsilon$-arguments
that appear there are the pattern for everything that follows.
